\section[Hecke 代数]{Hecke 代数（4 月 1 日）}

\srcpage{49}

上次我们除其他内容外讨论了上同调自守形式，例如 $H^*(G(\Z)\backslash G_\R/K)$，在幸运的情形它同构于 $G(\Z)$ 的群上同调。函数域类比更简单：离散集 $\Bun_G(C)(\F_q)$ 上的函数。

这里你可能会问一个重要问题：这应该与某个 $4$ 维 TFT $A_G$ 有关。这个理论给一个 $3$-流形 $Y$ 附上什么？当 $G$ 是环面时，我们把 $Y$ 上 $G$-联络空间的上同调附给 $Y$。但一般地不清楚该给 $3$-流形附什么：这是一个未决问题。这不会让「用拓扑场论看几何 Langlands」的计划注定失败：猜想本身是关于附给曲面的范畴的，而这些我们知道怎么定义。

回忆

\[\begin{aligned}
\Bun_G(C)(\F_q)&\cong G(F_C)\backslash G(\A_F)/G(\oo_{\A_F})\\
&=G(F_C)\backslash \prod_{x\in C}{}^{\prime} G(K_x)\Big/ \prod_{x\in C} G(\oo_x),
\end{aligned}\]

\srcpage{50}

其中限制乘积取遍 $C$ 的点 $x$，而 $C$ 是 $\F_q$ 上的曲线。这些就是我们的自守形式。顺着阿贝尔情形的思路，我们需要某个东西的作用来对角化并做谱分解；但上述双陪集空间不是群，所以不能立刻看出该怎么办。

但我们确实有形如 $X/K$ 的空间，带有群 $G$ 作用在 $X$ 上（这个 $G$ 未必是之前的 $G$），其中 $K\subset G$。所以一般的问题是：当商来自更大的对称群时，什么作用在商上的函数上？在上面的描述中，我们有巨大的对称性——每个点的 Laurent 级数群；但我们只商掉了其中一部分，即每个点的 Taylor 级数。若右边不取商，会有

\[G(F_C)\backslash \prod_{x\in C}{}^{\prime} G(K_x),\]

而它有 $\prod_x G(K_x)$ 的作用。这是许多份 $G(K_x)$，它们彼此交换！不过每个单独的群非交换。当我们再商掉 $\prod' G(\oo_x)$，这个对称性仍然存在。这是 Hecke 代数对称的一例；今天剩下的时间，我们在几个例子中讨论 Hecke 代数。

作为一个例子，设有限群 $G$ 作用在集合 $X$ 上，$K\subset G$，取 $\C[X/K]$。陪集空间 $X/K$ 有 $N_G(K)/K$ 的作用，但这并不总是很有用：对某些 $G,K$ 的选择它是平凡的。不过 $X/K$ 上的函数有更好的作用。

我们知道代数 $\C[G]$ 作用在向量空间 $\C[X]$ 上，且 $\C[X/K]=\C[X]^K$，即 $K$-不变量。那么什么作用在 $G$-表示的 $K$-不变量上？取 $K$-不变量是函子 $(-)^K:\Rep_G\to\Vect$，它可典范地表出：

\[(-)^K=\Hom_{\C[G]}(V_{G,K},-).\]

这是因为作为 $K$-表示，$V^K=\Hom_{\C[K]}(\C,V|_K)$，再用 Hom-张量伴随

\[\Hom_{\C[K]}(\C,V|_K)=\Hom_{\C[G]}(\C[G]\otimes_{\C[K]}\C,V).\]

这称为 Frobenius 互反律。\footnote{$\C[G]\otimes_{\C[K]}V$ 常称为 $V$ 的诱导表示，但一般地这是余诱导表示；对有限群二者一致，这会在后面造成混淆。}由于 $G$ 有限，可以把 $G$ 上的函数与测度等同：

\[V_{G,K}=\C[G]\otimes_{\C[K]}\C=\{f:G\to\C\mid f(gh)=f(g)\}.\]

所以这些是 $\C[G/K]=\C[G]^K$ 上的函数。当 $G$ 无限时，得到的将是测度而不是函数。

\begin{definition}
Hecke 代数 $\cH_{G,K}$ 是 $K$-不变量函子的自同态代数。
\end{definition}

由于取 $K$-不变量可表出，

\begin{align*}
\cH_{G,K}&=\End_{\Rep_G}V_{G,K}=\Hom_{\C[G]}(V_{G,K},V_{G,K})\\
&=(V_{G,K})^K=\C[G/K]^K=\C[K\backslash G/K].
\end{align*}

大部分时候，这就是双陪集上的函数出现在数学中的原因。先验地，双陪集上的函数没有代数结构，但这个构造给出了典范的代数结构。

一个叫 Barr–Beck 定理的抽象废话允许我们做如下论证：$K$-不变量函子 $\Rep_G\to\Vect$ 经过从 $\Rep_G$ 到 $\cH_{G,K}$-模范畴的函子分解：$K$-不变量典范地带有这个额外结构。\yizhu{Barr–Beck 定理即伴随函子的单子性／Beck 定理；原文称其为「抽象废话」是课堂戏称。}{A（原文 p.50）}这个遗忘函子忘掉的东西比想象中少：在 $\Rep_G$ 中考虑这样的满子范畴，其 $K$-不变量「足够大」，意思是它们作为 $G$-表示由自己的 $K$-不变量生成。例如对不可约表示，这只是要求存在某个非零的 $K$-不变向量。为了好玩，我们称这些表示为球面表示（spherical representation）。

结论是：从球面表示到 $\cH_{G,K}$-模范畴的遗忘函子是范畴等价！这里没有什么特别表示论的东西；全都是关于伴随函子的同义反复。换一种说法：球面范畴由 $V_{G,K}$ 生成。

其中大部分在更大的一般性下成立，例外是我们几次用到「有限集上的函数与测度等同」。但 Barr–Beck 定理的应用仍然有效。

\begin{example}
球谐函数（spherical harmonics）理论取 $\mathrm{SO}_3$ 在 $S^2$ 上的作用并线性化，考虑 $\mathrm{SO}_3$ 在 $L^2(S^2)$ 上的作用。这引向高中化学课上见过的 $s,p,d,f$ 轨道图像。
\end{example}

\srcpage{51}

我们可以把 $S^2$ 写成陪集空间 $\mathrm{SO}_3/\mathrm{SO}_2$：$G=\mathrm{SO}_3$，$K=\mathrm{SO}_2$。于是 $G$ 的表示是球面的，若它作为 $L^2(S^2)$ 的直和项出现，这等价于它有 $\mathrm{SO}_2$-不变量。所有 $\mathrm{SO}_3$-表示都是球面的；但若换成 $\mathrm{SU}_2$，会发现只有一半是。这就是「球面」一词的来历。

下面更具体地描述 $\C[K\backslash G/K]$ 上的代数结构。对 $\C[G]$，我们有推拉描述：

\begin{tikzcd}
 & G\times G \arrow[dl, "\pi_1"'] \arrow[dr, "\pi_2"] \arrow[d, "\mu"] & \\
 G & G & G
\end{tikzcd}

$f,g\in\C[G]$ 的乘积由

\[f\boxtimes h=\pi_1^*f\cdot \pi_2^*h\]

给出，要定义乘法（卷积），就用 $\mu$ 前推。关键在于，这一切对第一个 $G$ 因子上的左 $K$-作用与第二个 $G$ 拷贝上的右 $K$-作用都不变：

\begin{tikzcd}
 & K\backslash G\times G/K \arrow[dl, "\pi_1"'] \arrow[dr, "\pi_2"] \arrow[d, "\mu"] & \\
 K\backslash G & K\backslash G/K & G/K
\end{tikzcd}

在 $K\backslash G\times G/K$ 上还有一份「中间」的 $K$ 作用，所以可以取「平衡积」，令 $(g\cdot k,g')=(g,k\cdot g')$，记作 $-\times_K-$。而这个图表对模掉这份额外的 $K$ 不变：

\begin{tikzcd}
 & K\backslash G\times_K G/K \arrow[dl, "\pi_1"'] \arrow[dr, "\pi_2"] \arrow[d, "\mu"] & \\
 K\backslash G/K & K\backslash G/K & K\backslash G/K
\end{tikzcd}

乘法相同：卷积。可以把它想成子代数 $\C[G]^{K\times K}\subset\C[G]$（其中一个 $K$ 作用在右边）。其某种版本可以推广到层等等。

双陪集空间的另一种想法是

\[K\backslash G/K\cong (G/K\times G/K)/G,\]

其中 $G$ 对角作用。因此

\[\C[K\backslash G/K]=\C[G/K\times G/K]^G,\]\yizhu{原文此处印作 $\C[G\times K\times G/K]^G$，与上下文（它落在 $\C[G/K\times G/K]$ 中）及数学都不符，应为 $\C[G/K\times G/K]^G$；译文按后者排。}{A（原文 p.51）}

它落在 $\C[G/K\times G/K]$ 中。对任意集合 $X$，$X\times X$ 上的函数有典范的代数结构，即矩阵或 $\End(\C[X])$。矩阵乘法也有推拉描述

\begin{tikzcd}
 & X\times X\times X \arrow[dl, "\pi_{12}"'] \arrow[dr, "\pi_{23}"] \arrow[d, "\pi_{13}"] & \\
 X\times X & X\times X & X\times X
\end{tikzcd}

\srcpage{52}

即沿外侧两个投影拉回、相乘，再前推。这给出理解卷积下群代数的另一种方式：$G$-不变的 $G\times G$ 矩阵，其中 $G$ 对角作用。这正是上面那个描述的实质。

这个例子是所有 Hecke 代数背后的基本例子；你若遇到叫 Hecke 代数的东西，就该默想这个例子。

好，群胚时间。这本质上是用更抽象的语言说同一个例子，但对我们中更「叠」的人可能有帮助。群胚 $pt/G$ 是 $\Bun_G(pt)$：只有一个主 $G$-丛，它有 $G$ 个自同构。这可以用纤维积图表说明：

\begin{tikzcd}
 pt\times_{pt/G}pt \arrow[r] \arrow[d] & pt \arrow[d] \\
 pt \arrow[r] & pt/G
\end{tikzcd}

现在选子群 $K\subset G$。映射 $K\to G$ 给出 $pt/K\to pt/G$，而拉回

\begin{tikzcd}
 {} \arrow[r] \arrow[d] & pt/K \arrow[d, "G/K"] \\
 pt/K \arrow[r, "G/K"'] & pt/G
\end{tikzcd}

存在，且恰是 $K\backslash G/K$。即

\[K\backslash G/K=pt/K\times_{pt/G}pt/K.\]

一般地，给定映射 $\pi:X\to Y$（TODO：这里必须满射吗？），$X\times_Y X$ 上的函数可以想成分块对角矩阵。类似的描述给出 $\C[K\backslash G/K]$ 上的代数结构。这意味着 Hecke 代数是群胚代数（groupoid algebra）的一例（群代数的推广）。

指定空间 $X$ 上的 $G$-作用，等价于知道叠的映射 $X/G\to pt/G$。（若叠不是你的菜，可以把 $X$ 当成集合，$X/G$ 与 $pt/G$ 都是群胚。）然后可以由拉回

\begin{tikzcd}
 X \arrow[r] \arrow[d] & X/G \arrow[d] \\
 pt \arrow[r] & pt/G
\end{tikzcd}

找回 $X$。

所以回忆我们问过：对 $G$-空间 $X$，什么作用在 $X/K$ 上？我们看到了函数情形的答案，也看到对群而言答案是 $N_G(K)/K$。但我们其实能看到群胚作用：有拉回图表

\begin{tikzcd}
 X/K \arrow[r] \arrow[d] & pt/K \arrow[d] \\
 X/G \arrow[r] & pt/G
\end{tikzcd}

即 $X/K\to pt/K$（等价于带 $K$-作用的 $X$，如上所见）是从 $X\to pt/G$ 拉回的。

\begin{tikzcd}
 & X\times_{pt/K}K\backslash G/K \arrow[dl] \arrow[dr] \arrow[d] & \\
 X/K & X\times_K G/K & X/K
\end{tikzcd}

这展开来说就是：在 $X/K$ 之上有一个主 $G$-丛，我们在做相伴丛构造（TODO：我觉得这里漏了些东西）。

\begin{example}
若 $K=\{1\}$，则 $K\backslash G/K=G$，Hecke 代数是 $\C[G]$。
\end{example}

\srcpage{53}

\begin{example}
若 $K=G$，则 $K\backslash G/K$ 不平凡：我们有两个 $G$-作用和一个 $G$ 拷贝，所以得到 $pt/G$。但 Hecke 代数是平凡的：$\cH_{G,G}\cong\C$。
\end{example}

\begin{example}
取 Lie 群 $G$ 与 Borel 子群 $B$，令 $K=B(\F_q)\subset G(\F_q)$。例如 $G=\mathrm{GL}_n$ 时，$B$ 可取上三角矩阵子群。

此时 $\cH_{G,K}$ 是有限 Hecke 代数，即旗簇 $G(\F_q)/B(\F_q)$ 上 $B(\F_q)$-不变函数构成的代数。作为向量空间，Hecke 代数是 $\C[W]$，其中 $W$ 是 $G$ 的 Weyl 群（例如 $G=\mathrm{GL}_n$ 时为 $S_n$）；但代数结构不同——它依赖有限域的大小 $q$，某种意义上是以通常代数结构为底的 $\C[W]$ 的 $q$-形变。

记 $\cB=G(\F_q)/B(\F_q)$，把双陪集空间写成 $G(\F_q)\backslash \cB\times\cB$，这给出旗的相对版本。若 $G=\mathrm{SL}_2$，则 $\cB=\PP^1_{\F_q}$，于是用不依赖坐标的方式只能问两条线是否相等，所以应该预期 $B(\F_q)$ 在旗簇对 $\cB\times\cB$ 上的轨道只有两个元素：两旗相等的轨道，与两旗不等的轨道。这个预期是对的。因此 Hecke 代数是二维的，由 $\{1,T\}$ 张成；在相等轨道上 $T*T=q$，在不等轨道上 $T*T=q-1$。也就是说，你问有多少东西与原来的旗不同：若相同，有 $q$ 个点；若不同，只有 $q-1$ 个。注意 $q=\#(\A^1_{\F_q})$，$q-1=\#(\G_{m,\F_q})$。合起来，

\[T^2=(q-1)T+q,\]

或等价地 $(T+1)(T-q)=0$。若令 $q=1$，得 $T^2=1$，即 Weyl 群 $\Z/2$ 的群代数。这有点诡异（一元域是什么？），但提示了一般地这个 Hecke 代数是 $\C[\Z/2]$ 的 $q$-形变。按我们的构造，Hecke 代数是 $\C$-代数，但关系是 $q$ 的有限 Laurent 级数，所以 $\cH_q$ 其实是 $\Z[q,q^{-1}]$-代数，可以在 $q=1$ 取值。所以也许 $G(\F_1)=W$，尽管没有办法严格说清。

这展示了一个重要点：此情形 Hecke 代数是 Weyl 群的 $q$-形变，因此至少和 Weyl 群一样非交换。而 Weyl 群可以相当非交换，例如 $\mathrm{GL}_n$ 得到对称群 $S_n$ 的 $\C[S_n]$。Hecke 代数一般并不交换。
\end{example}

但在某些情形它们交换。对任意群 $K$，令 $G:=K\times K$，$K$ 对角嵌入 $G$。则 $\cH_{G,K}$ 交换！双陪集空间是

\[K\backslash (K\times K)/K=K/K,\]

其中 $K$ 通过共轭作用在 $K$ 上。取函数时得到的代数结构是类函数，即 $K/K$ 上的函数逐点相乘，这交换。事实上这是

\[\cH_{K\times K,K}=\End_{\C[K]\otimes\C[K]^{\mathrm{op}}}\C[K]=Z(\C[K]).\]

这是代数的中心，也叫 $\C[K]$ 的第零 Hochschild 同调。中心等同于 $\End_{A\otimes A^{\mathrm{op}}}(A)$ 并非群代数的特殊事实；对一般结合代数都成立。所以这个 Hecke 代数不仅交换，而且以非常稳健的方式交换。

我们将看到，作用在自守形式上的 Hecke 代数——它们很大很有趣，因为我们从巨大的对称性出发——是交换代数！这是 Langlands 纲领的诞生。事实上它们因结构性原因而交换，但这不那么显然。我们将用拓扑场论来证明。
